\chapter{Clinically Plausible Values}
\label{app:clinical-plausible-values}

This appendix lists clinically plausible (extreme–but–possible) minima and maxima for numeric features in the sepsis–AKI ICU cohort (first 24 h), intended as \emph{winsorization} bounds during preprocessing. Ranges reflect ICU/sepsis/AKI physiology rather than healthy outpatients and are aligned with the unit-harmonized guardrails used in the curated dataset.

\vspace{0.75\baselineskip}
\section*{Demographics}
\vspace{0.75\baselineskip}
\renewcommand{\arraystretch}{1.12}
\setlength{\tabcolsep}{6pt}
\begin{table}[ht!]\centering\small
\begin{tabular}{lrrc}
\toprule
\textbf{Variable} & \textbf{Min} & \textbf{Max} & \textbf{Units} \\
\midrule
Age & 18 & 120 & years \\
Admission weight & 30 & 300 & kg \\
\bottomrule
\end{tabular}
\end{table}

\noindent\textit{Notes.} Adult ICUs typically include patients $\ge$18\,y; extreme body weights spanning severe malnutrition to massive obesity motivate 30–300\,kg as a pragmatic span.

\vspace{0.75\baselineskip}
\section*{Timing anchors (descriptive only)}
\vspace{0.75\baselineskip}
\begin{table}[ht!]\centering\small
\begin{tabular}{lrrc}
\toprule
\textbf{Variable} & \textbf{Min} & \textbf{Max} & \textbf{Units} \\
\midrule
Time to death & 0 & 2{,}880 & h \\
Time to AKI onset & 0 & 168 & h \\
Time to sepsis onset & 0 & 72 & h \\
Time to mechanical ventilator start & 0 & 120 & h \\
Time to vasopressor start & 0 & 120 & h \\
\bottomrule
\end{tabular}
\end{table}

\noindent\textit{Notes.} Onset times cluster near ICU admission in sepsis and AKI; the upper bounds allow delayed recognition/therapy within the early ICU course.

\vspace{0.75\baselineskip}
\section*{Vital signs (first 24 h)}
\vspace{0.75\baselineskip}
\begin{table}[ht!]\centering\small
\begin{tabular}{lrrc}
\toprule
\textbf{Variable} & \textbf{Min} & \textbf{Max} & \textbf{Units} \\
\midrule
Heart rate & 20 & 250 & beats/min \\
Systolic blood pressure & 50 & 300 & mmHg \\
Diastolic blood pressure & 30 & 200 & mmHg \\
Mean arterial pressure & 30 & 200 & mmHg \\
Respiratory rate & 4 & 50 & breaths/min \\
Oxygen saturation & 50 & 100 & \% \\
\bottomrule
\end{tabular}
\end{table}

\noindent\textit{Notes.} Sustained tachyarrhythmias can transiently reach 200–250\,bpm (\citealt{ClevelandPSVT}); severe shock and hypertensive crisis motivate BP limits (\citealt{GSTTStages,AnesthesiaHypotension}). Profound hypoxemia may drive SpO$_2$ towards 50–60\% while values $>$100\% are non-physiologic (\citealt{MayoHypoxemia,FPNotebookO2Sat}).

\vspace{0.75\baselineskip}
\section*{Fluid balance and urine output (first 24 h)}
\vspace{0.75\baselineskip}
\begin{table}[ht!]\centering\small
\begin{tabular}{lrrc}
\toprule
\textbf{Variable} & \textbf{Min} & \textbf{Max} & \textbf{Units} \\
\midrule
Net fluid balance (24 h) & $-10{,}000$ & $+15{,}000$ & mL \\
Fluid balance per kg (24 h) & $-100$ & $+200$ & mL/kg \\
Urine output (24 h) & 0 & 20{,}000 & mL \\
Urine output per kg (24 h) & 0 & 200 & mL/kg \\
\bottomrule
\end{tabular}
\end{table}

\noindent\textit{Notes.} Bounds reflect aggressive resuscitation/diuresis and anuria–polyuria extremes early in critical illness.

\vspace{0.75\baselineskip}
\section*{Blood gases and oxygenation (first 24 h)}
\vspace{0.75\baselineskip}
\begin{table}[ht!]\centering\small
\begin{tabular}{lrrc}
\toprule
\textbf{Variable} & \textbf{Min} & \textbf{Max} & \textbf{Units} \\
\midrule
Arterial pH & 6.5 & 7.7 & --- \\
PaCO$_2$ & 10 & 100 & mmHg \\
PaO$_2$ & 30 & 500 & mmHg \\
PF ratio (PaO$_2$/FiO$_2$) & 0 & 600 & mmHg \\
\bottomrule
\end{tabular}
\end{table}

\noindent\textit{Notes.} Reported survivals at pH near 6.5–6.8 are exceptional (\citealt{pH6point68}); permissive hypercapnia and ventilatory failure bound PaCO$_2$ (\citealt{PermissiveHypercapnia,MedscapeRespAcidosis}). On FiO$_2$=1.0, PaO$_2$ can reach $\sim$400–500\,mmHg; PF ratios span 0–600\,mmHg (\citealt{TargetPO2,FPNotebookO2Sat}).

\vspace{0.75\baselineskip}
\section*{Electrolytes and acid–base (first 24 h)}
\vspace{0.75\baselineskip}
\begin{table}[ht!]\centering\small
\begin{tabular}{lrrc}
\toprule
\textbf{Variable} & \textbf{Min} & \textbf{Max} & \textbf{Units} \\
\midrule
Sodium & 100 & 190 & mmol/L \\
Potassium & 2.0 & 8.0 & mmol/L \\
Chloride & 70 & 130 & mmol/L \\
Bicarbonate & 5 & 40 & mmol/L \\
Anion gap & 3 & 30 & mmol/L \\
\bottomrule
\end{tabular}
\end{table}

\noindent\textit{Notes.} Bounds cover high–gap metabolic acidosis and severe alkalosis; see \citet{StatPearlsHAGMA,DKAGuideline} for context.

\vspace{0.75\baselineskip}
\section*{Renal function and glycemia (first 24 h)}
\vspace{0.75\baselineskip}
\begin{table}[ht!]\centering\small
\begin{tabular}{lrrc}
\toprule
\textbf{Variable} & \textbf{Min} & \textbf{Max} & \textbf{Units} \\
\midrule
Creatinine (serum) & 0.3 & 15.0 & mg/dL \\
Blood urea nitrogen & 5 & 200 & mg/dL \\
Glucose (plasma) & 20 & 800 & mg/dL \\
\bottomrule
\end{tabular}
\end{table}

\noindent\textit{Notes.} Severe AKI/ESRD without dialysis can manifest with creatinine 10–15\,mg/dL and BUN $>$100\,mg/dL; hyperosmolar states often show glucose 600–800+\,mg/dL (\citealt{ApolloCreatinine,MedNewsUremia,StatPearlsHHS}).

\vspace{0.75\baselineskip}
\section*{Calcium, magnesium, phosphate (first 24 h)}
\vspace{0.75\baselineskip}
\begin{table}[ht!]\centering\small
\begin{tabular}{lrrc}
\toprule
\textbf{Variable} & \textbf{Min} & \textbf{Max} & \textbf{Units} \\
\midrule
Calcium (total) & 4.0 & 15.0 & mg/dL \\
Magnesium (serum) & 0.5 & 6.0 & mg/dL \\
Phosphate (serum) & 0.5 & 15.0 & mg/dL \\
\bottomrule
\end{tabular}
\end{table}

\noindent\textit{Notes.} Life-threatening calcium extremes and clinically significant magnesium/phosphate derangements motivate these bounds (\citealt{MerckHypocalcemia,HypercalcemiaReview}).

\vspace{0.75\baselineskip}
\section*{Hepatic and biliary panel (first 24 h)}
\vspace{0.75\baselineskip}
\begin{table}[ht!]\centering\small
\begin{tabular}{lrrc}
\toprule
\textbf{Variable} & \textbf{Min} & \textbf{Max} & \textbf{Units} \\
\midrule
Alkaline phosphatase & 20 & 1{,}200 & U/L \\
AST & 5 & 10{,}000 & U/L \\
ALT & 5 & 10{,}000 & U/L \\
Total bilirubin & 0 & 40 & mg/dL \\
Direct bilirubin & 0 & 30 & mg/dL \\
Gamma–glutamyl transferase & 0 & 2{,}000 & U/L \\
Albumin (serum) & 1.0 & 6.0 & g/dL \\
\bottomrule
\end{tabular}
\end{table}

\noindent\textit{Notes.} Massive transaminase elevations in ischemic/toxic injury (\citealt{ALT1000,ALF_Acet}) and marked cholestatic patterns (\citealt{MayoALP}) justify the upper bounds; severe hypoalbuminemia is common in critical illness (\citealt{StatPearlsAlbumin}).

\vspace{0.75\baselineskip}
\section*{Hematology and coagulation (first 24 h)}
\vspace{0.75\baselineskip}
\begin{table}[ht!]\centering\small
\begin{tabular}{lrrc}
\toprule
\textbf{Variable} & \textbf{Min} & \textbf{Max} & \textbf{Units} \\
\midrule
Hemoglobin & 3 & 20 & g/dL \\
Platelet count & 5 & 1{,}000 & $\times 10^9$/L \\
INR & 0.8 & 10.0 & ratio \\
aPTT & 15 & 200 & s \\
White blood cell count & 0 & 200 & $\times 10^9$/L \\
\bottomrule
\end{tabular}
\end{table}

\noindent\textit{Notes.} Extremes of anemia/polycythemia (\citealt{EJCRIM_SevereAnemia,HighHbMerits}), critical thrombocytopenia and thrombocytosis (\citealt{StatPearlsPlateletTx,PMCThrombocythemia}), and profound coagulopathy (\citealt{INRHighBleeding}) set these bounds. Leukocytosis may rarely exceed 200$\times 10^9$/L.

\vspace{0.75\baselineskip}
\section*{Inflammation and tissue injury (first 24 h)}
\vspace{0.75\baselineskip}
\begin{table}[ht!]\centering\small
\begin{tabular}{lrrc}
\toprule
\textbf{Variable} & \textbf{Min} & \textbf{Max} & \textbf{Units} \\
\midrule
C–reactive protein & 0 & 500 & mg/L \\
Lactate dehydrogenase & 50 & 10{,}000 & U/L \\
Troponin (cardiac) & 0 & 100 & ng/mL \\
BNP or NT–proBNP & 0 & 100{,}000 & pg/mL \\
Lactate (serum) & 0.5 & 20 & mmol/L \\
\bottomrule
\end{tabular}
\end{table}

\noindent\textit{Notes.} Severe sepsis often elevates lactate; levels $\gtrsim$4–10\,mmol/L indicate high risk (\citealt{AAFPSepsisLactate}). Very high troponin reflects extensive myocardial injury.

\vspace{0.75\baselineskip}
\section*{Severity scores (day 1)}
\vspace{0.75\baselineskip}
\begin{table}[ht!]\centering\small
\begin{tabular}{lrrc}
\toprule
\textbf{Variable} & \textbf{Min} & \textbf{Max} & \textbf{Units} \\
\midrule
APACHE III score & 0 & 299 & points \\
SOFA total (24 h) & 0 & 24 & points \\
SOFA subscore (each organ) & 0 & 4 & points \\
\bottomrule
\end{tabular}
\end{table}

\noindent\textit{Notes.} Ranges follow score definitions (\citealt{APACHEIII,PhysiopediaSOFA}).

\vspace{0.75\baselineskip}
\paragraph{Implementation.} Bounds are applied after unit harmonization and basic plausibility checks; outliers beyond these limits are capped, not dropped. These limits match the clinical guardrails used to compute laboratory QC fractions in the main dataset.
